\documentclass[letterpaper, 11pt]{article}
\usepackage[utf8]{inputenc}
\usepackage[spanish]{babel}
\usepackage[T1]{fontenc}
\usepackage{charter}
\usepackage{csquotes}
\usepackage[margin=3cm]{geometry}
\usepackage{fancyhdr}
\usepackage{minted}
\usepackage{hyperref}
\usemintedstyle{vs}

% --- DEFINICIÓN DE ESTILOs DE PÁGINA ---
\fancypagestyle{sections}{
    \fancyhf{} 
    %\fancyhead[L]{Proyecto Final LTII} --> Texto a la izquierda
    \fancyhead[R]{Planteamiento de secciones}
    \fancyfoot[C]{\thepage}
    \setlength{\headheight}{15pt}
    \renewcommand{\headrulewidth}{0.4pt}
}
% -----------------------------


\title{Proyecto Final: Laboratorio Temático II}
\author{@Char[50] \textit{Carlos T. Mendoza}}
\date{\today}

\renewcommand{\thesection}{\Alph{section}}
\renewcommand{\thesubsection}{\roman{subsection}}

\begin{document}
\maketitle
\tableofcontents

\newpage
\section{Antecedentes}
    Steam ha existido desde \textit{septiembre de 2003} y aunque en sus inicios tuvo un arranque algo \enquote{atropellado}, hoy en día es la plataforma más enorme y poderosa en el ámbito del \textit{Gaming} para \textit{PC}. Incluso ahora \enquote{y gracias al desarollo Proton \footnote{Capa intermedia que traduce a Linux las instrucciones de un videojuego nativo de Windows.}} ha incursionado en el mundo de los \enquote{PC's de mano} con su \textbf{SteamDeck} y, a últimas fechas, en el de \enquote{consolas} con la mítica \textbf{SteamMachine}. 

    Más allá de lo asombroso que pueda ser, es inegable que ostenta un monopolio extremadamente sobrado \textit{\enquote{controlando al rededor de un 70 a 75\% de la distribución de videojuegos en formato digital para PC}}. Por lo anterior, mi intención con este proyecto no es sólo apoyarme de sus datos para ilustrar algunos tópicos interesantes, sino también poner en contexto al espectador tanto de lo \enquote{bueno} como lo \enquote{malo} de la plataforma y su influencia. 

    Los datos utilizados para este proyecto fueron seleccionados tanto por su fuente \enquote{puesto que vienen directamente de la plataforma en cuestion} como por su relevancia en el ámbito comercial y del \textit{Gaming}. 


\section{Contenido}
    \subsection{Sección A} 

    \subsection{Sección B} 

    \subsection{Sección C} 
    

\newpage
\section{Conclusión}


\end{document}






